\subsection{Interaction of an alkali atom with a magnetic field}
We must now consider the effect of a weak external magnetic field on the energy levels of an
alkali atom. This will produce the \textbf{Zeeman Effect}, and will result in further splitting
of the energy levels. What is meant by ``weak'' magnetic field? If the resulting splitting is
very small compared to the Hyperfine Splitting (HFS), the magnetic field is said to be weak.
This will be the case in all the experiments discussed here.
 
The magnetic field $\mathbf{B}$ causes $\mathbf{F}$ to precess about it at the Larmour
frequency, with $M$ being the component of $\mathbf{F}$ along the field direction.
The Hamiltonian that accounts for the interaction of the electronic and nuclear magnetic
moments with the external field is :
\begin{equation}
    H = ha\, \mathbf{I} \cdot \mathbf{J}
      - \frac{\mu_J}{J}\,\mathbf{J} \cdot \mathbf{B}
      - \frac{\mu_I}{I}\,\mathbf{I} \cdot \mathbf{B}
    \label{eq:hamiltonian_Bfield}
\end{equation}
where $\mu_J$ is the total electronic magnetic dipole moment (spin coupled to orbit) and
$\mu_I$ is the nuclear magnetic dipole moment. The magnetic field splits each $F$ level into
$2F + 1$ sublevels that are approximately equally spaced.
 
\subsubsection*{The Land\'{e} $g$-Factor}
 
Neglecting the nucleus, the magnetic interaction energy is written via the vector model as :
\begin{equation}
    W = \frac{(\mathbf{L} + 2\mathbf{S})\cdot\mathbf{J}}{J(J+1)}\,\mu_0 MB
      = g_J\,\mu_0 MB
    \label{eq:mag_energy}
\end{equation}
where the Land\'{e} $g$-factor $g_J$ evaluates to :
\begin{equation}
    g_J = 1 + \frac{J(J+1) + S(S+1) - L(L+1)}{2J(J+1)}
    \label{eq:gJ}
\end{equation}
In terms of this $g$-factor, the interaction energy with a magnetic field is:
\begin{equation}
    W = -g_J\,\mu_0\,BM
    \label{eq:W_electron}
\end{equation}
where $B$ is the magnitude of the magnetic field and $M$ is the projection of the electron
spin along $\mathbf{B}$. For rubidium, $J = S = \tfrac{1}{2}$, giving $g_J = 2$
(the measured value is $g_J = 2.00232$).
 
\subsubsection*{Hyperfine Zeeman Effect}
 
When the nuclear interaction is included, the $g$-factor becomes :
\begin{equation}
    g_F = g_J\,\frac{F(F+1) + J(J+1) - I(I+1)}{2F(F+1)}
    \label{eq:gF}
\end{equation}
and the interaction energy is :
\begin{equation}
    W = g_F\,\mu_0\,BM
    \label{eq:W_hyperfine}
\end{equation}
where the direct interaction of the nuclear moment with the magnetic field is neglected.
 
\subsubsection*{The Breit-Rabi Equation}
 
For the purposes of this experiment, terms quadratic in the magnetic field must be
considered. Diagonalizing Eq.~(\ref{eq:hamiltonian_Bfield}) by perturbation theory yields
the \textbf{Breit-Rabi equation} :
\begin{equation}
    W(F,M) = -\frac{\Delta W}{2(2I+1)} - \mu_I BM
    \pm \frac{\Delta W}{2}
    \left(1 + \frac{4Mx}{2I+1} + x^2\right)^{1/2}
    \label{eq:breit_rabi}
\end{equation}
where
\begin{equation}
    \mu_I = -g_I\frac{\mu_0}{I},
    \qquad
    x = \frac{(g_J - g_I)\,\mu_0\,B}{\Delta W}
    \label{eq:BR_params}
\end{equation}
and $\Delta W$ is the hyperfine energy splitting.
 
\subsubsection*{Regions of the Breit-Rabi Diagram}
 
The Breit-Rabi diagram can be divided into three main regions :
\begin{enumerate}
    \item \textbf{Zeeman region} ($x \approx 0$) : Energy splittings vary \emph{linearly}
          with the applied field; $M$ is a good quantum number.
    \item \textbf{Paschen-Back region} ($x > 2$) : Energy levels are again linear in $B$,
          corresponding to full decoupling of $\mathbf{I}$ and $\mathbf{J}$; $m_I$ and
          $m_J$ become good quantum numbers.
    \item \textbf{Intermediate field region} : Energy levels are nonlinear in $B$;
          $\mathbf{I}$ and $\mathbf{J}$ are partially decoupled and $M$ is no longer a
          good quantum number.
\end{enumerate}
At all fields the relation $M = m_I + m_J$ holds. In the optical pumping experiment we are
concerned with small magnetic fields, where the levels are either linear or exhibit only a
small quadratic dependence on $B$.